\chapter{PSDL: The Python-Semantic Deppy Language}
\label{ch:psdl}

\PSDL{} is the strong proof language at the heart of the
\Deppy{} verification pipeline.  It is what fills the
\texttt{psdl\_proof:} field of a \texttt{verify} block.  This chapter
is a complete reference: design rationale, syntax, statement-level
mapping, soundness arguments, and worked examples.

The implementation is in
\texttt{deppy/proofs/psdl.py} (1553 lines as of writing).  We cite
specific line ranges as we go.

\section{Design intent}
\label{sec:psdl-design}

The unifying design choice of \PSDL{} is that \emph{the proof script
is literal Python}.  \texttt{ast.parse} reads it; the
\texttt{\_Compiler} class walks the AST and maps each statement to:
\begin{enumerate}[leftmargin=1.7em]
  \item a \Deppy{} \texttt{ProofTerm} that
    \texttt{ProofKernel.verify} type-checks;
  \item a line of \Lean{}~4 tactic-mode text emitted into the
    enclosing \texttt{by ...} block.
\end{enumerate}

Two design constraints follow:

\paragraph{1. No new syntax.} Users should not have to learn a
ninth proof language.  \PSDL{} \emph{is} Python --- if it parses as
Python, \PSDL{} can in principle accept it.  The semantics of each
construct, however, is \emph{tactic-language} semantics, not Python
runtime semantics: \texttt{if isinstance(other, Point)} does not
\emph{run} an \texttt{isinstance} check, it \emph{introduces a
fibration}.

\paragraph{2. Two-target compilation.} Every statement must produce
both a \texttt{ProofTerm} and a \Lean{} tactic.  When a statement
cannot be compiled to one of the two, the result records an error:
\begin{lstlisting}
@dataclass
class PSDLError:
    line_no: int
    line_text: str
    message: str

@dataclass
class PSDLCompileResult:
    proof_term: Optional[ProofTerm] = None
    lean_tactic: str = ""
    term_repr: str = ""
    errors: list[PSDLError] = field(default_factory=list)
    hints: list[str] = field(default_factory=list)
    show_goal_requested: bool = False
    foundations_referenced: list[str] = field(default_factory=list)
\end{lstlisting}
The \texttt{ok} property is \texttt{True} iff both the
\texttt{ProofTerm} is constructed and \texttt{errors} is empty.
This is the gate.

\section{The two-vocabulary structure}
\label{sec:two-vocab}

\PSDL's vocabulary splits cleanly into two halves, kept separate
to keep clear what is borrowed from \CIC{} and what is intrinsic to
\Deppy.

\paragraph{\CIC{} core.}  $\Pi$, $\Sigma$, identity types, induction
principles, substitution, equivalences, congruence, function
extensionality, transport, path induction.  Same expressive power as
\Lean{} / \Coq{} / \Agda{}.

\paragraph{Cubical / duck / effect Python primitives.}  Things that
make sense only because the underlying language being reasoned about
is Python:
\begin{itemize}[leftmargin=1.7em]
  \item \texttt{DuckPath} dispatch on shared method sets;
  \item $\Diamond$ / $\Box$ effect modalities for value/exception
    fibres;
  \item \texttt{isinstance} as a fibration;
  \item Kan-filling in cubical control-flow analyses;
  \item cocycle constructors for inter-procedural cohomology.
\end{itemize}

The lookup tables in this chapter group the constructors by which
half of the vocabulary they inhabit.

\section{Surface mapping: control flow}
\label{sec:psdl-control}

\subsection{\texttt{if isinstance(x, T):}}

A standard idiom in Python: \texttt{if isinstance(x, T): \dots else:
\dots}.  In \PSDL{} this is the canonical fibration construction.

\begin{lstlisting}[language=psdl]
if isinstance(x, T):
    <value-fibre body>
else:
    <complement-fibre body>
\end{lstlisting}

The compilation, from \texttt{psdl.py}'s \texttt{\_compile\_if}
(lines 609--642):

\begin{itemize}[leftmargin=1.7em]
  \item Detects the \texttt{isinstance} test pattern via
    \texttt{\_is\_isinstance\_test}.
  \item Builds a \texttt{Fiber(scrutinee=Var(x), type\_branches=[(T,
    body), (\_, else\_branch)], exhaustive=True)}.
  \item Emits \Lean: \verb|by_cases h : isinstance x T|, then the
    \texttt{·} bullets for each fibre.
\end{itemize}

When the condition is \emph{not} \texttt{isinstance}, the compiler
falls through to a generic \texttt{Cases} on the test expression:

\begin{lstlisting}[language=psdl]
if x > 0:
    ...
else:
    ...
\end{lstlisting}

becomes a \texttt{Cases(scrutinee=Var("x > 0"), branches=[("then",
\dots), ("else", \dots)])} and \Lean{} \verb|by_cases h : x > 0|.

\subsection{\texttt{try / except}}

\begin{lstlisting}[language=psdl]
try:
    <value-fibre body>
except TypeError:
    <handler body>
\end{lstlisting}

\texttt{\_compile\_try} (lines 645--666) folds:
\begin{enumerate}[leftmargin=1.7em]
  \item the \texttt{try}-body into an
    \texttt{EffectWitness(effect="value", proof=$P_v$)};
  \item each \texttt{except} handler into a
    \texttt{ConditionalEffectWitness(precondition="raises Exc",
    effect="exception", proof=$P_e$, target=Exc)}.
\end{enumerate}
The \Lean{} side emits \texttt{-- try block (effect modality)} as
a comment and the handler body inline.  The kernel verifies the
inner proofs; the modal layer is recorded as effect-fibre witnesses.

\subsection{\texttt{for} and \texttt{while}}

\begin{lstlisting}[language=psdl]
for x in xs:
    <step body>
\end{lstlisting}

compiles via \texttt{\_compile\_for} (lines 668--685) to
\texttt{ListInduction(var=x, nil\_proof=Refl(\dots),
cons\_proof=step\_body)}.  The \Lean{} emits
\verb|induction xs with | followed by \texttt{| nil => rfl} and
\texttt{| cons \_ \_ ih =>}.

\begin{lstlisting}[language=psdl]
while n:
    <step body>
\end{lstlisting}

compiles to \texttt{NatInduction(var=n, base\_proof=Refl,
step\_proof=body)}.  The default base case is \texttt{rfl}; users
who want a richer base case use the explicit
\texttt{induction(n, base=\dots, step=\dots)} call form.

\subsection{\texttt{match} / \texttt{case}}

Python 3.10's structural pattern-matching:

\begin{lstlisting}[language=psdl]
match x:
    case A:
        <body_A>
    case B:
        <body_B>
\end{lstlisting}

compiles to \texttt{Cases(scrutinee=Var("x"), branches=[(A, $P_A$),
(B, $P_B$)])} and \Lean{} \verb|match x with | \verb|| A => ... | B => ...|.

\section{Surface mapping: logical / equality tactics}
\label{sec:psdl-logical}

\subsection{\texttt{assert P} and discharge modes}

\begin{lstlisting}[language=psdl]
assert P                  # default: discharge by Z3
assert P, "z3"            # explicit Z3 discharge
assert P, "by F"          # cite foundation F
assert P, "by lemma L"    # cite lemma L
\end{lstlisting}

The compiler (\texttt{\_compile\_assert}, lines 710--777) reads the
optional message argument as a discharge selector:
\begin{itemize}[leftmargin=1.7em]
  \item \texttt{"z3"} (or no message) $\to$ \texttt{Cut(\dots,
    Z3Proof(formula=P, binders=\{\}), \dots)}; \Lean{} emits
    \verb|have h : P := by omega -- z3-discharged|.
  \item \texttt{"by F"} where $F$ is in the foundations dict $\to$
    \texttt{Cut(\dots, AxiomInvocation(F, "foundations"), \dots)};
    \Lean{} emits \verb|have h : P := Foundation_F|.
  \item \texttt{"by lemma L"} $\to$ \texttt{Cut(\dots,
    AxiomInvocation(L, ""), \dots)}; \Lean{} emits
    \verb|have h : P := L|.
\end{itemize}

The \texttt{Cut} bound name is auto-generated as \texttt{h\_<lineno>}.

\subsection{\texttt{apply(F)} and bare-name citation}

\begin{lstlisting}[language=psdl]
apply(F)
F                # alias: bare-name reference
\end{lstlisting}

Both compile to \texttt{AxiomInvocation(name=F, module=\dots)}.  When
$F$ is a registered foundation, the module is \texttt{"foundations"}
and the \Lean{} emits \verb|exact Foundation_F|.  Otherwise the
module is empty and \Lean{} emits \verb|exact F|.

\subsection{\texttt{inline}, \texttt{unfold}, \texttt{rewrite}}

\begin{lstlisting}[language=psdl]
inline(method, expected_expr)    # body translation assertion
unfold(method)                   # plain Unfold
method.unfold()                  # alias

rewrite(L)                       # rewrite using L
L.rewrite()                      # alias
\end{lstlisting}

\texttt{inline(method, expected\_expr)} emits an \texttt{Unfold}
proof term whose body is \texttt{Refl(Var(expected\_expr))} ---
i.e., it asserts that after unfolding \texttt{method}, the body
matches \texttt{expected\_expr} up to reflexivity.  The \Lean{} side
emits \verb|unfold method  -- expected: <expected_expr>|.

\texttt{unfold} alone emits \texttt{Unfold(method, Refl(...))} with
no explicit RHS.  \texttt{rewrite} emits \texttt{Rewrite(lemma,
Refl(...))} and \Lean{} \verb|rw [lemma]|.

\subsection{\texttt{return}, \texttt{pass}, \texttt{raise}}

\begin{lstlisting}[language=psdl]
return expr      # Refl
pass             # Structural("trivial"); Lean: trivial
raise TypeError                 # exception fibre marker
raise NotImplementedError       # vacuous (gated Structural)
\end{lstlisting}

\texttt{return expr} compiles to \texttt{Refl(Var(safe(expr)))}; the
\Lean{} prints the comment \verb|-- return <expr>| followed by
\verb|rfl|.

\texttt{raise NotImplementedError} is the canonical
\emph{vacuous-branch} marker: emit
\texttt{Structural("vacuous: NotImplementedError")} and \Lean{}
\verb|exact absurd ‹_› (by simp)|.  The soundness gate counts this
\texttt{Structural} leaf and the certificate's
\texttt{certify\_or\_die} accepts it only when the verification
context (e.g.\ the exception fibre of an \texttt{if isinstance})
makes the branch genuinely unreachable.

\subsection{Local lemmas and assignments}

\begin{lstlisting}[language=psdl]
h: T = expr                       # annotated: typed local lemma
L = lemma(name, statement, by)    # explicit lemma constructor
\end{lstlisting}

The annotated assignment compiles via \texttt{\_compile\_ann\_assign}
(lines 808--827) to a \texttt{Cut} whose
\texttt{hyp\_type} is \texttt{RefinementType(base\_type=PyBoolType,
predicate=T)} and whose \texttt{lemma\_proof} is the compiled RHS.
The \Lean{} side emits \verb|have h : T := <RHS>|.

The \texttt{lemma(...)} form (compiled in \texttt{\_compile\_func\_call}
under \texttt{lemma}, lines 1020--1040) is an explicit
constructor: \texttt{lemma(name, statement, by)} where \texttt{by} is
the discharge text.

\section{Surface mapping: CIC core}
\label{sec:psdl-cic}

The \CIC{} core surface is what gives \PSDL{} its theoretical
expressive power.  Every constructor in this section emits a real
\Lean{} tactic that elaborates against the corresponding \CIC{} rule.

\subsection{Quantifiers and induction}

\begin{lstlisting}[language=psdl]
forall(x, T, body_proof)       # ∀ introduction
exists(x, T, witness, body)    # ∃ introduction (with explicit witness)

induction(n, base=..., step=...)   # explicit Nat induction
\end{lstlisting}

\texttt{forall(x, T, body)} emits \verb|intro (x : T)| and recurses
into \texttt{body}.  The \texttt{ProofTerm} is a \texttt{Cut} whose
\texttt{hyp\_type} is the refinement \texttt{\{x : T\}} and whose
\texttt{lemma\_proof} is the de Bruijn assumption
\texttt{Assume(name=x)}.

\texttt{exists(x, T, witness, body)} emits
\verb|refine ⟨witness, ?_⟩|.  The \texttt{ProofTerm} is a
\texttt{Cut} whose \texttt{lemma\_proof} is \texttt{Refl(Var(witness))}
and whose \texttt{body\_proof} continues with the supplied body.

\texttt{induction(n, base=$P_0$, step=$P_s$)} emits
\verb|induction n with| followed by the base/step alternatives, and
the term is \texttt{NatInduction(var=n, base\_proof=$P_0$,
step\_proof=$P_s$)}.

\subsection{Equality and substitution}

\begin{lstlisting}[language=psdl]
refl(x)                       # Refl[x]
sym(eq)                       # Eq.symm
trans(eq1, eq2)               # Eq.trans
cong(f, eq)                   # congrArg f eq
subst(eq, motive, body)       # rw [eq]
rew(eq, body)                 # bare rewrite
\end{lstlisting}

These compile to the obvious kernel terms:
\texttt{Refl}, \texttt{Sym(p)}, \texttt{Trans(p, q)}, \texttt{Cong(f,
p)}, \texttt{Rewrite(eq, body)}.  \Lean{} emits the corresponding
tactics: \verb|exact (eq).symm|, \verb|exact (eq1).trans (eq2)|,
\verb|exact congrArg f eq|, \verb|rw [eq]|.

\subsection{Path constructors}

\begin{lstlisting}[language=psdl]
path_concat(p, q)             # composite path → PathComp
path_inv(p)                   # inverse → Sym
path_ap(f, p)                 # functorial action → Ap
J(motive, refl_case, x, y, p) # path induction → TransportProof
\end{lstlisting}

\texttt{path\_concat} compiles to \texttt{PathComp(p, q)} and
\Lean{} \verb|exact (p).trans (q)|.  \texttt{path\_inv} is
\texttt{Sym}.  \texttt{path\_ap(f, p)} is \texttt{Ap(f, p)} and
\Lean{} \verb|exact congrArg f p|.

The Martin-L\"of $J$ rule is approximated as a transport: the
implementation (lines 1192--1209 of \texttt{psdl.py}) builds a
\texttt{TransportProof} whose type-family is the motive,
path-proof is an \texttt{AxiomInvocation} on the supplied path, and
base-proof is the compiled refl-case.  The \Lean{} side emits the
explicit \verb|exact @Eq.rec _ _ <motive> <refl> _ <p>|.  This is a
faithful encoding for first-order motives; higher-dimensional
motives are accepted but not specially decomposed.

\subsection{Equivalence and univalence}

\begin{lstlisting}[language=psdl]
iso(forward, backward, l_inv, r_inv)   # construct an Equiv
equiv_to_path(e)                       # ua-style transport
\end{lstlisting}

\texttt{iso} compiles to a \texttt{Cut} whose hypothesis is
\texttt{Equiv-of-iso}.  The \Lean{} emits
\verb|exact Equiv.mk fwd bwd l_inv r_inv|.

\texttt{equiv\_to\_path(e)} is the bridge to univalence: it builds
a \texttt{TransportProof} with type-family
\texttt{Var("\_motive")}, path-proof
\texttt{AxiomInvocation(name="ua", module="psdl\_local")}, and base
proof the compiled $e$.  The \Lean{} side emits a comment ---
\Lean{} core does not have univalence; the user must add it through
\texttt{lean\_imports} if needed.

\section{Surface mapping: cubical / duck / effect}
\label{sec:psdl-cubical-duck}

This section is what distinguishes \PSDL{} from a literal \CIC{}
tactic language.  These constructors model the
intra-Python phenomena that the \Deppy{} metatheory was built to
handle.

\subsection{Duck dispatch}

\begin{lstlisting}[language=psdl]
dispatch(receiver, "method")
\end{lstlisting}

Compiled (lines 1249--1262) to \texttt{DuckPath(base\_type=PyObj,
target\_type=PyObj)}.  The \Lean{} emits
\verb|exact DuckPath_resolve recv "method"|.  In a flat-Int
preamble the \texttt{DuckPath\_resolve} is admitted as
\verb|axiom|; in a richer preamble where method tables are
declared, it elaborates against an actual Σ-projection.

\subsection{Kan filling}

\begin{lstlisting}[language=psdl]
kan_fill(open_box, filling)
\end{lstlisting}

Compiles (lines 1263--1272) to a \texttt{Patch} term wrapping the
filling proof and emits \verb|apply Kan.fill|.  The kernel verifies
the wrapped proof; the cubical-coherence check is performed by
\texttt{KanFill} (when constructed directly) or skipped for the
sugar form.  See \Cref{ch:cohomology} for when the full
\texttt{KanFill} is engaged.

\subsection{Cocycle construction}

\begin{lstlisting}[language=psdl]
cocycle(level, witness)
\end{lstlisting}

Builds a \texttt{Cocycle(level, label="cocycle\_<lvl>",
proof=Refl(Var(witness)))} and emits a comment in the \Lean{}.
The cocycle's full structure (boundary pairs, overlap proofs,
boundary-zero witness) is not exposed in the sugar form; users who
need it construct \texttt{Cocycle} directly via the kernel API.

\subsection{Fibration split}

\begin{lstlisting}[language=psdl]
fibre_split(scrutinee, [(T1, body1), (T2, body2)])
\end{lstlisting}

The explicit form of \texttt{if isinstance}: a \texttt{Fiber} with
the listed branches.  The \Lean{} side emits a comment
\verb|-- explicit fibration on <scrut>|.  Used when the user wants
multiple type fibres rather than just \texttt{T} vs.\ \texttt{not T}.

\subsection{Diamond / Box modalities}

\begin{lstlisting}[language=psdl]
pure(v)        # □-introduction
bind(m, k)     # ◇ Kleisli composition
expect(m, fallback)   # exception-fibre handler
lift(v)        # ◇-introduction
lower(m, witness)     # ◇-elimination with purity witness
\end{lstlisting}

These are the surface forms for the modal layer.  In the kernel:
\begin{itemize}[leftmargin=1.7em]
  \item \texttt{lift(v)} $\to$ \texttt{EffectWitness(effect="value", proof=$P_v$)}.
  \item \texttt{bind(m, k)} $\to$ a \texttt{Cut} whose lemma proof is
    $m$ and body is $k$.
  \item \texttt{expect(m, fallback)} $\to$ a chain of effect
    witnesses ending in a \texttt{ConditionalEffectWitness}.
  \item \texttt{lower(m, witness)} $\to$ the inner proof of $m$
    (assumes the purity witness is supplied at the same trust
    level).
  \item \texttt{pure} alone is structural (the marker).
\end{itemize}

\subsection{Homotopy and squares}

\begin{lstlisting}[language=psdl]
homotopy(p, q)                # path between paths
square(top, bot, left, right) # 2-square diagram
\end{lstlisting}

\texttt{homotopy(p, q)} compiles to \texttt{PathComp(p, q)};
\texttt{square(\dots)} chains four paths via two compositions.
There is no native ``Square'' \texttt{ProofTerm}; the user obtains
the semantically richer form by constructing
\texttt{KanFill(dimension=2, faces=\dots)} directly.

\section{Top-level user declarations}
\label{sec:psdl-toplevel}

\PSDL{} is a Python module, which means the user can declare
inductives, axioms, theorems at top level.

\subsection{\texttt{@inductive class T:}}

\begin{lstlisting}[language=psdl]
@inductive
class Tree:
    Leaf: int
    Node: (Tree, Tree)
\end{lstlisting}

The \texttt{\_compile\_class\_def} method (lines 481--536) walks the
class body looking for \texttt{AnnAssign} (named constructors with
explicit types) or \texttt{FunctionDef} (constructor-as-method).
The \Lean{} side emits an \verb|inductive| declaration:
\begin{lstlisting}[language=lean4]
inductive Tree where
  | Leaf : int
  | Node : Tree → Tree → Tree
\end{lstlisting}

The kernel records this as a structural ``inductive Tree (2 ctors)''
that can be cited inside the same \PSDL{} block via
\texttt{Tree.\_\_dict\_\_}-style references.

\subsection{\texttt{@axiom def} and \texttt{@theorem def}}

\begin{lstlisting}[language=psdl]
@axiom
def magic() -> Path[Int, 0, 1]: ...

@theorem
def my_lemma(x: Nat) -> Path[Nat, x + 0, x]:
    induction(x) on Nat:
        case Z(): refl(0)
        case S(k, ih): cong(succ, ih)
\end{lstlisting}

\texttt{@axiom} produces a Lean \verb|axiom magic : Path Int 0 1|
and registers it in the \PSDL{} compiler's \texttt{foundations}
dict so subsequent \texttt{apply(magic)} resolves correctly.

\texttt{@theorem} elaborates the body as a \PSDL{} proof and emits a
real \verb|theorem my_lemma : <type> := by <body>|.  The kernel
records a \texttt{Cut} whose hypothesis name is the theorem name.

The \texttt{@protocol} decorator declares a structural protocol
similar to PEP 544's \texttt{Protocol} class.

\section{Compositional structure: how multiple statements become one ProofTerm}
\label{sec:psdl-composition}

A \PSDL{} block is a \emph{list} of statements.  The compiler chains
them via \texttt{Cut}: the first statement's proof becomes the
\texttt{lemma\_proof} of a \texttt{Cut} whose body is the recursion
on the rest.

\begin{lstlisting}[caption={\texttt{compile\_block} (lines 405--431),
the chaining engine.}]
def compile_block(self, stmts, indent=0):
    if not stmts:
        return Structural("empty proof block")
    if len(stmts) == 1:
        return self._compile_stmt(stmts[0], indent)
    head, *rest = stmts
    head_term = self._compile_stmt(head, indent)
    rest_term = self.compile_block(rest, indent)
    if isinstance(head_term, Structural) and head_term.description in (
            "show_goal", "show_term", "trivial", "pure"):
        return rest_term
    return Cut(
        hyp_name=f"_psdl_step_{getattr(head, 'lineno', 0)}",
        hyp_type=PyObjType(),
        lemma_proof=head_term,
        body_proof=rest_term,
    )
\end{lstlisting}

The exception in the middle is the optimisation: pure-noise
statements (\verb|show_goal()|, \verb|hint("...")|, \verb|pure()|)
do \emph{not} introduce a real \texttt{Cut} but pass through to the
next statement.  This keeps the proof tree compact.

\section{Worked compilation: \texttt{dist\_nonneg\_psdl}}
\label{sec:psdl-worked}

We trace through the \texttt{dist\_nonneg\_psdl} block from
\texttt{sympy\_geometry.deppy} statement-by-statement.

The \PSDL{} source is:

\begin{lstlisting}[language=psdl]
if isinstance(other, Point):
    inline(Point.distance, sqrt(sum_zip_sub_sq(self, other)))
    apply(Real_sqrt_nonneg)
    assert sum_zip_sub_sq(self, other) >= 0, "z3"
else:
    raise NotImplementedError
\end{lstlisting}

The compiler walks this top-down:

\paragraph{1. The outer \texttt{if isinstance} statement.}
\texttt{\_compile\_if} fires.  It detects the \texttt{isinstance}
shape via \texttt{\_is\_isinstance\_test}, which returns
\texttt{("other", "Point")}.  The compiler emits the \Lean{} lines:

\begin{lstlisting}[language=lean4]
by_cases h : isinstance other Point
· -- value-fibre: isinstance(other, Point)
  ...
· -- complement-fibre: ¬ isinstance(other, Point)
  ...
\end{lstlisting}

It recursively compiles the body of the if-branch (a 3-statement
block) and the body of the else-branch (a 1-statement block), then
returns:

\begin{verbatim}
Fiber(
  scrutinee=Var("other"),
  type_branches=[
    (PyObjType(), <if-body>),
    (PyObjType(), <else-body>),
  ],
  exhaustive=True,
)
\end{verbatim}

\paragraph{2. The if-body, statement 1: \texttt{inline(Point.distance, sqrt(...))}.}
\texttt{\_compile\_func\_call} dispatches on \texttt{fname=="inline"}
(lines 951--963).  \texttt{target = "Point.distance"};
\texttt{expected = "sqrt(sum\_zip\_sub\_sq(self, other))"}.

The compiled term is
\texttt{Unfold(func\_name="Point.distance",
proof=Refl(Var("sqrt\_sum\_zip\_sub\_sq\_self\_other")))} (with the
\texttt{\_safe} identifier transform applied to the expected
expression).  The \Lean{} emits
\verb|unfold Point.distance  -- expected: sqrt(...)|.

\paragraph{3. The if-body, statement 2: \texttt{apply(Real\_sqrt\_nonneg)}.}
\texttt{\_compile\_func\_call} dispatches on \texttt{fname=="apply"}
(lines 932--941).  \texttt{target = "Real\_sqrt\_nonneg"}, which is
in the \texttt{foundations} dict (registered by the sidecar parser).

Term: \texttt{AxiomInvocation(name="Real\_sqrt\_nonneg",
module="foundations")}.  The compiler appends
\texttt{"Real\_sqrt\_nonneg"} to \texttt{foundations\_referenced}.
\Lean{}: \verb|exact Foundation_Real_sqrt_nonneg|.

\paragraph{4. The if-body, statement 3: \texttt{assert ... , "z3"}.}
\texttt{\_compile\_assert} (lines 710--777) reads \texttt{stmt.msg ==
ast.Constant("z3")}, sets mode to \texttt{"z3"}, builds
\texttt{Z3Proof(formula="sum\_zip\_sub\_sq(self, other) >= 0",
binders=\{\})} and emits
\verb|have h_<lineno> : sum_zip_sub_sq(self, other) >= 0 := by omega|.

The full term is
\begin{verbatim}
Cut(
  hyp_name="h_<lineno>",
  hyp_type=RefinementType(
      base_type=PyBoolType(),
      predicate="sum_zip_sub_sq(self, other) >= 0",
  ),
  lemma_proof=Z3Proof(...),
  body_proof=Refl(Var("_assert_body")),
)
\end{verbatim}

\paragraph{5. Composition of the if-body.}
\texttt{compile\_block} chains the three statements through nested
\texttt{Cut}s, with \texttt{Refl(Var("\_assert\_body"))} as the
ultimate body.

\paragraph{6. The else-body: \texttt{raise NotImplementedError}.}
\texttt{\_compile\_raise} (lines 789--805) detects
\texttt{NotImplementedError} and returns
\texttt{Structural("vacuous: NotImplementedError")} with \Lean{}
\verb|exact absurd ‹_› (by simp)|.

\paragraph{7. The full ProofTerm.}
Putting it all together (with some abbreviation):

\begin{verbatim}
Fiber(scrutinee=Var("other"),
      type_branches=[
        (PyObjType(),
          Cut(hyp_name="_step_1",
            lemma_proof=Unfold("Point.distance", Refl(...)),
            body_proof=Cut(hyp_name="_step_2",
              lemma_proof=AxiomInvocation("Real_sqrt_nonneg",
                                          "foundations"),
              body_proof=Cut(hyp_name="h_<lineno>",
                hyp_type=RefinementType(predicate="sum_zip... >= 0"),
                lemma_proof=Z3Proof("sum_zip... >= 0"),
                body_proof=Refl(Var("_assert_body")))))),
        (PyObjType(), Structural("vacuous: NotImplementedError")),
      ])
\end{verbatim}

\paragraph{8. The \Lean{} block.}

\begin{lstlisting}[language=lean4]
by
  by_cases h : isinstance other Point
  · -- value-fibre: isinstance(other, Point)
    unfold Point.distance  -- expected: sqrt(sum_zip_sub_sq(self, other))
    exact Foundation_Real_sqrt_nonneg
    have h_<lineno> : sum_zip_sub_sq(self, other) >= 0 := by omega
      -- z3-discharged
  · -- complement-fibre: ¬ isinstance(other, Point)
    -- vacuous: NotImplementedError
    exact absurd ‹_› (by simp)
\end{lstlisting}

\paragraph{9. Trust composition.}
The trust on this \texttt{ProofTerm} is
$\min(\mathtt{KERNEL}, \mathtt{AXIOM\_TRUSTED}, \mathtt{Z3\_VERIFIED}, \mathtt{STRUCTURAL\_CHAIN}) = \mathtt{STRUCTURAL\_CHAIN}$.  \texttt{certify\_or\_die} permits this
because the \texttt{Structural} leaf is in a vacuous-branch position
(the exception fibre), and it counts as a non-failure.

\section{Soundness arguments per constructor}
\label{sec:psdl-soundness}

For each compiler entry point we sketch the soundness story.

\paragraph{If-isinstance.}
\texttt{Fiber(scrutinee, type\_branches, exhaustive=True)} is
sound when:
\begin{enumerate}[leftmargin=1.7em]
  \item every \texttt{type\_branch}'s proof verifies under the local
    refinement \texttt{isinstance(scrutinee, T)};
  \item the type-branches \emph{cover} the dynamic type of the
    scrutinee.
\end{enumerate}
The PSDL form generates \texttt{exhaustive=bool(stmt.orelse)}: True
when the user wrote an \texttt{else}.  When the user did \emph{not}
provide an \texttt{else}, the kernel adds an implicit ``other''
branch with \texttt{Structural("no else for isinstance(...)")}, which
downgrades trust to \texttt{STRUCTURAL\_CHAIN}.

\paragraph{Try-except.}
The chain \texttt{EffectWitness(value)} +
\texttt{ConditionalEffectWitness(precond="raises Exc", \dots)} is
sound only if \emph{both} fibres are addressed.  The PSDL form
respects this: each \texttt{except} clause produces a corresponding
witness; the value fibre is built from the \texttt{try}-body.  We do
\emph{not} assert that \texttt{except} catches every possible
exception: if the body can raise \texttt{ValueError} but the user
only wrote \texttt{except TypeError}, the resulting modal type
admits the \texttt{ValueError} fibre as part of the residual
$\Diamond$ structure.

\paragraph{Z3 assert.}
\texttt{Z3Proof(formula, binders)} is sound when SMT returns
\texttt{unsat} on the negation in the typed-Z3 encoder.  Trust
\texttt{Z3\_VERIFIED}; we admit \Zthree's answer.

\paragraph{Apply.}
\texttt{AxiomInvocation(name, module)} is sound when the named
axiom's statement matches the goal up to the kernel's
\texttt{PatternMatcher}.  Trust \texttt{AXIOM\_TRUSTED}.  The
PSDL form additionally records the axiom in
\texttt{foundations\_referenced} so the certificate's provenance is
honest about which foundations / axioms a verify block depends on.

\paragraph{Inline / unfold / rewrite.}
\texttt{Unfold(name, proof)} is sound when the named definition
unfolds to a term equivalent to the proof's claim.  In the current
implementation, this check is structural rather than reduction-based;
\Lean's elaborator performs the actual definitional unfolding when
the certificate is round-tripped.

\paragraph{Vacuous \texttt{raise}.}
\texttt{Structural("vacuous: <descr>")} is permitted only in
positions where the \emph{branch is unreachable}: typically inside
the exception-fibre face of a \texttt{Fiber} or
\texttt{ConditionalEffectWitness}.  The soundness gate
(\Cref{ch:soundness}) flags these and \texttt{certify\_or\_die}
accepts them only when no counterexample is found that contradicts
the unreachability.

\section{Limitations and open work}
\label{sec:psdl-limitations}

\PSDL{} as currently implemented:
\begin{itemize}[leftmargin=1.7em]
  \item Does not support universe polymorphism.  All types live in
    \texttt{Type 0}.
  \item Does not support refinement types beyond a single
    \texttt{predicate} string field.  The kernel does not check the
    predicate; the \Lean{} side does, by elaborating the predicate.
  \item Does not support higher inductive types.  The
    \texttt{@inductive} decorator generates ordinary inductives.
  \item Does not support coinductive specs.
  \item The escape hatch \texttt{lean("\dots raw \dots")} compiles to
    a \texttt{LeanProof} term and emits the raw text, but \Deppy{}'s
    kernel does \emph{not} type-check the raw fragment ---
    \Lean{} does.
\end{itemize}

These gaps are documented in \texttt{AUDIT.md} as items B (type-system
expressiveness), F (advanced cubical/homotopy), and G (module
structure).  The escape hatch \texttt{lean(\dots)} suffices in
practice for any property a user can prove directly in \Lean.

\section{Summary}
\label{sec:ch6-summary}

\PSDL{} is \emph{the} surface for \Deppy{}'s strong proof language:
\begin{itemize}[leftmargin=1.7em]
  \item Source is literal Python.  No new syntax.
  \item Each statement compiles to a \texttt{ProofTerm} the kernel
    verifies and a \Lean{} tactic the elaborator verifies.
  \item Vocabulary splits into a \CIC{} core (Π/Σ/induction/transport)
    and Python-specific primitives (fibrations, $\Diamond$/$\Box$,
    duck dispatch, Kan filling, cocycles).
  \item The compiler's structure is direct: an AST walker with one
    \texttt{\_compile\_<X>} method per AST node kind.
  \item Composition is via \texttt{Cut}-chaining; pure-noise
    statements pass through unchanged.
  \item Soundness follows from the kernel's verification rules
    (\Cref{ch:cubical}); trust composes by \texttt{min\_trust}.
\end{itemize}

The next chapter follows the pipeline from start to finish: how
\texttt{run\_sidecar} ingests a \texttt{.deppy} file and produces a
\Lean{} certificate.
